% Options for packages loaded elsewhere
\PassOptionsToPackage{unicode}{hyperref}
\PassOptionsToPackage{hyphens}{url}
\PassOptionsToPackage{dvipsnames,svgnames,x11names}{xcolor}
%
\documentclass[
]{article}
\usepackage{amsmath,amssymb}
\usepackage{iftex}
\ifPDFTeX
  \usepackage[T1]{fontenc}
  \usepackage[utf8]{inputenc}
  \usepackage{textcomp} % provide euro and other symbols
\else % if luatex or xetex
  \usepackage{unicode-math} % this also loads fontspec
  \defaultfontfeatures{Scale=MatchLowercase}
  \defaultfontfeatures[\rmfamily]{Ligatures=TeX,Scale=1}
\fi
\usepackage{lmodern}
\ifPDFTeX\else
  % xetex/luatex font selection
  \setmainfont[]{Inter}
\fi
% Use upquote if available, for straight quotes in verbatim environments
\IfFileExists{upquote.sty}{\usepackage{upquote}}{}
\IfFileExists{microtype.sty}{% use microtype if available
  \usepackage[]{microtype}
  \UseMicrotypeSet[protrusion]{basicmath} % disable protrusion for tt fonts
}{}
\makeatletter
\@ifundefined{KOMAClassName}{% if non-KOMA class
  \IfFileExists{parskip.sty}{%
    \usepackage{parskip}
  }{% else
    \setlength{\parindent}{0pt}
    \setlength{\parskip}{6pt plus 2pt minus 1pt}}
}{% if KOMA class
  \KOMAoptions{parskip=half}}
\makeatother
\usepackage{xcolor}
\usepackage[margin=18mm]{geometry}
\usepackage{longtable,booktabs,array}
\usepackage{calc} % for calculating minipage widths
% Correct order of tables after \paragraph or \subparagraph
\usepackage{etoolbox}
\makeatletter
\patchcmd\longtable{\par}{\if@noskipsec\mbox{}\fi\par}{}{}
\makeatother
% Allow footnotes in longtable head/foot
\IfFileExists{footnotehyper.sty}{\usepackage{footnotehyper}}{\usepackage{footnote}}
\makesavenoteenv{longtable}
\setlength{\emergencystretch}{3em} % prevent overfull lines
\providecommand{\tightlist}{%
  \setlength{\itemsep}{0pt}\setlength{\parskip}{0pt}}
\setcounter{secnumdepth}{-\maxdimen} % remove section numbering
\ifLuaTeX
\usepackage[bidi=basic]{babel}
\else
\usepackage[bidi=default]{babel}
\fi
\babelprovide[main,import]{spanish}
\ifPDFTeX
\else
\babelfont{rm}[]{Inter}
\fi
% get rid of language-specific shorthands (see #6817):
\let\LanguageShortHands\languageshorthands
\def\languageshorthands#1{}
\ifLuaTeX
  \usepackage{selnolig}  % disable illegal ligatures
\fi
\IfFileExists{bookmark.sty}{\usepackage{bookmark}}{\usepackage{hyperref}}
\IfFileExists{xurl.sty}{\usepackage{xurl}}{} % add URL line breaks if available
\urlstyle{same}
\hypersetup{
  pdflang={es-CL},
  colorlinks=true,
  linkcolor={Maroon},
  filecolor={Maroon},
  citecolor={Blue},
  urlcolor={Blue},
  pdfcreator={LaTeX via pandoc}}

\title{Anexos contractuales — Plataforma CICUC}
\author{Investigación, Tecnología y Gestión ADYAC SpA}
\date{31 de julio de 2026}

\begin{document}
\maketitle

{
\hypersetup{linkcolor=}
\setcounter{tocdepth}{3}
\tableofcontents
}
\hypertarget{anexo-a--alcance-hitos-aceptaciuxf3n-y-soporte}{%
\section{Anexo A --- Alcance, hitos, aceptación y
soporte}\label{anexo-a--alcance-hitos-aceptaciuxf3n-y-soporte}}

\textbf{Cotización:} ADYAC-COT-CICUC-2026-001\\
\textbf{Versión:} 1.0 --- 31 de julio de 2026\\
\textbf{Estado:} Borrador sujeto a revisión y firma

Este anexo forma parte integral e inseparable de la cotización.

\hypertarget{1-objeto-y-luxedmites}{%
\subsection{1. Objeto y límites}\label{1-objeto-y-luxedmites}}

Adyac diseñará, desarrollará, validará y pondrá en marcha una plataforma
web para la gestión administrativa de estudios clínicos oncológicos de
CICUC. La plataforma apoya organización y preselección manual; no
determina elegibilidad, no recomienda tratamientos y no reemplaza el
protocolo ni el juicio profesional.

Quedan fuera del alcance inicial IA, extracción automática, voz, ficha
clínica, mensajería externa y elegibilidad automática.

\hypertarget{2-luxednea-base}{%
\subsection{2. Línea base}\label{2-luxednea-base}}

Durante el primer mes las partes aprobarán por escrito una línea base
con historias incluidas y excluidas, actores, permisos, flujos, datos,
requisitos no funcionales, criterios de aceptación, pruebas,
dependencias y responsables. No se tratarán datos reales ni se
desarrollarán funciones dependientes de decisiones clínicas o legales
bloqueantes antes de su aprobación.

\hypertarget{3-hitos-y-facturaciuxf3n}{%
\subsection{3. Hitos y facturación}\label{3-hitos-y-facturaciuxf3n}}

\begin{longtable}[]{@{}lll@{}}
\toprule\noalign{}
Hito & Entregables mínimos & Condición de facturación \\
\midrule\noalign{}
\endhead
\bottomrule\noalign{}
\endlastfoot
Mes 1 --- Descubrimiento & Línea base, arquitectura, modelo conceptual,
permisos provisionales, prototipo y ambiente demo & Aprobación escrita
de línea base y entrega \\
Mes 2 --- Núcleo & Autorización, inventario, ciclo de estudios,
versiones y criterios manuales & Aceptación escrita y corrección de
defectos bloqueantes \\
Mes 3 --- Operación & Flujos operativos acordados, auditoría/reportes
aprobados, documentación, capacitación y despliegue & Aceptación final y
entrega de código, accesos y documentación \\
\end{longtable}

Cada hito tiene un valor de \textbf{CLP 3.000.000 netos}, más IVA.
Total: \textbf{CLP 10.710.000 IVA incluido}. Cada cuota se factura
únicamente después de la aceptación escrita del hito y se paga dentro de
15 días corridos desde la recepción conforme del documento tributario.
Una extensión requiere acuerdo escrito sobre alcance, plazo y precio; no
se presume gratuita.

\hypertarget{4-aceptaciuxf3n}{%
\subsection{4. Aceptación}\label{4-aceptaciuxf3n}}

\begin{enumerate}
\def\labelenumi{\arabic{enumi}.}
\tightlist
\item
  Adyac notificará la entrega y acompañará evidencia.
\item
  CICUC tendrá 5 días hábiles para revisar contra criterios acordados.
\item
  La aceptación deberá constar por escrito; el silencio no constituye
  aceptación.
\item
  El rechazo identificará criterios incumplidos y evidencia
  reproducible.
\item
  Adyac corregirá sin costo defectos atribuibles al alcance y volverá a
  presentar.
\item
  Requisitos nuevos seguirán control de cambios.
\end{enumerate}

\hypertarget{5-dependencias-de-cicuc}{%
\subsection{5. Dependencias de CICUC}\label{5-dependencias-de-cicuc}}

CICUC designará una contraparte, entregará decisiones oportunamente,
proporcionará solo datos sintéticos o autorizados y validará roles,
privacidad, custodia, infraestructura y KPI. Los retrasos atribuibles a
estas dependencias ajustarán el calendario de común acuerdo y por
escrito.

\hypertarget{6-soporte-y-garantuxeda}{%
\subsection{6. Soporte y garantía}\label{6-soporte-y-garantuxeda}}

Se incluyen 12 meses desde la aceptación final: 30 días de
estabilización, garantía correctiva de 90 días, atención hábil de lunes
a viernes 09:00--18:00 Chile y hasta 6 horas mensuales no acumulables.
Objetivos de primera respuesta: crítica 4 horas hábiles; alta 8 horas;
media 1 día; baja 2 días. Nuevas funciones, rediseños, integraciones,
licencias, nube y operación 24/7 se cotizan aparte.

\hypertarget{7-tuxe9rmino-y-responsabilidad}{%
\subsection{7. Término y
responsabilidad}\label{7-tuxe9rmino-y-responsabilidad}}

Cualquiera de las partes podrá terminar por incumplimiento material no
subsanado dentro de 10 días hábiles desde notificación. CICUC pagará
entregables aceptados; Adyac entregará artefactos pagados y cooperará
razonablemente con la transición.

Cada parte responderá por daños directos acreditados que le sean
imputables. Se propone como límite general el monto neto pagado o
pagadero, salvo dolo, culpa grave u obligaciones irrenunciables.
Confidencialidad, datos personales e infracción de propiedad intelectual
deben acordarse en el contrato definitivo.

\hypertarget{8-facultad}{%
\subsection{8. Facultad}\label{8-facultad}}

La aceptación solo será válida por personas con facultades suficientes.
Toda orden de compra deberá identificar la cotización y los Anexos A, B
y C, con versión y fecha.

\textbf{Referencia:} Código Civil de Chile, especialmente artículos 1545
y 1546: \url{https://www.bcn.cl/leychile/navegar?idNorma=172986}.

\hypertarget{anexo-b--acuerdo-propuesto-de-tratamiento-de-datos}{%
\section{Anexo B --- Acuerdo propuesto de tratamiento de
datos}\label{anexo-b--acuerdo-propuesto-de-tratamiento-de-datos}}

\textbf{Cotización:} ADYAC-COT-CICUC-2026-001 · \textbf{Versión 1.0 ---
31 de julio de 2026}\\
\textbf{Estado:} Borrador para revisión jurídica, de privacidad,
seguridad y ética UC

\hypertarget{1-roles-y-finalidad}{%
\subsection{1. Roles y finalidad}\label{1-roles-y-finalidad}}

La Pontificia Universidad Católica de Chile/CICUC actuará como
responsable y Adyac únicamente como encargado, conforme a instrucciones
documentadas. La finalidad será prestar desarrollo, soporte y operación
técnica de la plataforma. Se prohíben publicidad, perfiles comerciales,
venta, cesión, usos propios y entrenamiento de modelos con datos de
CICUC.

\hypertarget{2-datos-y-titulares}{%
\subsection{2. Datos y titulares}\label{2-datos-y-titulares}}

Las categorías definitivas deberán aprobarse antes de producción.
Podrían incluir datos identificatorios mínimos, datos clínicos
sensibles, datos operacionales, usuarios y auditoría; los titulares
podrían ser candidatos, participantes y profesionales. Durante
desarrollo y pruebas se usarán exclusivamente datos sintéticos o
correctamente anonimizados.

\hypertarget{3-instrucciones-y-confidencialidad}{%
\subsection{3. Instrucciones y
confidencialidad}\label{3-instrucciones-y-confidencialidad}}

Adyac tratará datos solo según instrucciones; limitará acceso a personal
autorizado y sujeto a confidencialidad; registrará accesos relevantes;
notificará instrucciones potencialmente ilícitas y no revelará
información salvo autorización o mandato legal.

\hypertarget{4-seguridad-muxednima}{%
\subsection{4. Seguridad mínima}\label{4-seguridad-muxednima}}

\begin{itemize}
\tightlist
\item
  cifrado en tránsito y, cuando corresponda, en reposo;
\item
  mínimo privilegio, autenticación individual y segregación de
  ambientes;
\item
  auditoría, gestión de vulnerabilidades y secretos fuera del código;
\item
  respaldos, restauración, eliminación segura y control de
  exportaciones;
\item
  procedimiento probado de respuesta ante incidentes.
\end{itemize}

\hypertarget{5-infraestructura-y-subencargados}{%
\subsection{5. Infraestructura y
subencargados}\label{5-infraestructura-y-subencargados}}

Este borrador no autoriza una ubicación productiva ni proveedor
definitivo. Adyac informará previamente proveedor, servicios, países de
tratamiento, subencargados, medidas de seguridad y mecanismo para
transferencias internacionales. Nuevos subencargados seguirán el
mecanismo de aprobación UC.

\hypertarget{6-incidentes}{%
\subsection{6. Incidentes}\label{6-incidentes}}

Adyac notificará a CICUC sin dilación indebida y, como objetivo
contractual, dentro de 24 horas desde que tome conocimiento. Informará
naturaleza, datos y titulares potencialmente afectados, riesgos y
medidas; preservará evidencia y cooperará con la investigación. La
notificación no implica reconocimiento automático de responsabilidad.

\hypertarget{7-derechos-auditoruxeda-y-evaluaciones}{%
\subsection{7. Derechos, auditoría y
evaluaciones}\label{7-derechos-auditoruxeda-y-evaluaciones}}

Adyac asistirá razonablemente en derechos de titulares, evaluaciones de
impacto, consultas regulatorias y auditorías acordadas. No responderá
directamente a titulares salvo instrucción o exigencia legal.

\hypertarget{8-tuxe9rmino}{%
\subsection{8. Término}\label{8-tuxe9rmino}}

Al finalizar, Adyac devolverá los datos en formato acordado y eliminará
copias bajo su control, salvo obligación legal, emitiendo confirmación
escrita. Los respaldos se eliminarán conforme al ciclo documentado y
seguirán protegidos.

\hypertarget{9-marco-normativo}{%
\subsection{9. Marco normativo}\label{9-marco-normativo}}

Las partes considerarán la Ley 19.628, la Ley 20.584 y las exigencias de
la Ley 21.719 desde el 1 de diciembre de 2026. Este anexo debe adecuarse
a políticas UC y ser aprobado antes de utilizar datos personales reales.

\textbf{Referencias oficiales:}
\href{https://www.bcn.cl/leychile/Navegar?idNorma=141599}{Ley 19.628},
\href{https://www.bcn.cl/leychile/Navegar?idNorma=1039348}{Ley 20.584} y
\href{https://www.bcn.cl/leychile/navegar?idNorma=1209272}{Ley 21.719}.

\hypertarget{anexo-c--propiedad-intelectual-seguridad-y-continuidad}{%
\section{Anexo C --- Propiedad intelectual, seguridad y
continuidad}\label{anexo-c--propiedad-intelectual-seguridad-y-continuidad}}

\textbf{Cotización:} ADYAC-COT-CICUC-2026-001 · \textbf{Versión 1.0 ---
31 de julio de 2026}\\
\textbf{Estado:} Borrador sujeto a revisión y firma

\hypertarget{1-entregables-especuxedficos}{%
\subsection{1. Entregables
específicos}\label{1-entregables-especuxedficos}}

Pagados íntegramente los hitos aceptados, CICUC será titular de los
derechos patrimoniales que correspondan sobre desarrollos creados
específicamente por encargo: código fuente, migraciones, configuración
específica, documentación, pruebas, diseños y contenidos específicos
pagados. Recibirá repositorios, historial disponible, instrucciones de
build, dependencias y accesos bajo su control. Podrá usar, modificar y
encargar mantenimiento a terceros para fines institucionales.

\hypertarget{2-material-preexistente-de-adyac}{%
\subsection{2. Material preexistente de
Adyac}\label{2-material-preexistente-de-adyac}}

Adyac conserva herramientas, bibliotecas, plantillas, métodos y
componentes genéricos preexistentes o desarrollados independientemente.
Todo componente incorporado deberá identificarse previamente en un
registro. Sobre componentes indispensables, Adyac concede licencia
perpetua, irrevocable, mundial y libre de regalías, transferible junto
al sistema o a mantenedores de CICUC.

No se incluyen componentes de machine learning en el alcance inicial.
Cualquier incorporación futura exige un anexo específico sobre modelos,
datos, evaluación y derechos.

\hypertarget{3-terceros-y-datos}{%
\subsection{3. Terceros y datos}\label{3-terceros-y-datos}}

Adyac mantendrá inventario de dependencias y licencias y no incorporará
licencias incompatibles sin aprobación. CICUC conserva todos los
derechos sobre sus datos, protocolos, marcas y documentos; Adyac no
podrá reutilizarlos ni entrenar modelos.

\hypertarget{4-seguridad-de-entrega}{%
\subsection{4. Seguridad de entrega}\label{4-seguridad-de-entrega}}

Antes de puesta en marcha se entregará evidencia de pruebas,
autorización negativa, gestión de secretos, dependencias revisadas,
respaldo/restauración, logs sin información sensible, riesgos conocidos
y procedimiento de incidentes.

\hypertarget{5-continuidad}{%
\subsection{5. Continuidad}\label{5-continuidad}}

Ante término o cambio de proveedor, Adyac entregará artefactos pagados,
documentación y exportaciones acordadas, y cooperará razonablemente. No
podrá retener datos de CICUC como garantía de pago.

\hypertarget{6-registro-previo-a-la-firma}{%
\subsection{6. Registro previo a la
firma}\label{6-registro-previo-a-la-firma}}

\begin{longtable}[]{@{}ll@{}}
\toprule\noalign{}
Registro & Estado \\
\midrule\noalign{}
\endhead
\bottomrule\noalign{}
\endlastfoot
Componentes preexistentes de Adyac & Pendiente \\
Dependencias open source y licencias & Pendiente \\
Proveedor cloud, servicios y ubicación & Pendiente \\
Subencargados autorizados & Pendiente \\
Repositorios y administradores institucionales & Pendiente \\
Procedimiento de salida y formato de exportación & Pendiente \\
\end{longtable}

\textbf{Referencia oficial:} Ley 17.336 sobre Propiedad Intelectual,
artículo 8 y disposiciones relacionadas:
\url{https://www.bcn.cl/leychile/navegar?idNorma=28933}.

\end{document}
